% !TEX root = tarasov-qasp-zh.tex

\vspace*{1cm}

\chapter[俄文版编者的话]{俄文版编者的话}
\label{preface}

可以断言，还没有一个正在准备高校物理入学考试的学生，翻开过与本书类似的书。作者采用极其生动的对话形式，全面讨论了教学大纲中几乎所有的课题，特别是那些通常被认为难以理解的问题。书中对考生在物理入学考试中常犯的错误作了详细分析。学生会发现这是一本异常清晰而有趣的读物：它从多个角度处理复杂问题，并配有大量有助于更深入理解相关概念与思想的精彩插图。作者是莫斯科电子技术学院（Moscow Institute of Electronics Engineering）的教师，非常了解报考高校学生的总体训练水平，并有多年主持入学考试的经验。作者的专业学识，加上生动明晰的讲解，使本书成为准备物理考试的学生极为有用的辅导读物。\\
\begin{flushright}
\emph{G. 叶皮法诺夫}（G.~Epifanov）教授，物理数学博士
\end{flushright}

\chapter{作者前言}
\label{foreword}

本书是为准备高校物理入学考试的学生编写的辅助读物，采用作者（\textbf{教师：}）与好问的读者（\textbf{学生：}）对话的形式。这对于分析考生在入学考试中的常见错误、复习同一问题的多种解法、讨论物理理论中的疑难问题，都特别方便。

书中讨论了大量中学物理的问题与习题，此外还配备了供课后自测的习题（附答案）。大部分问题与习题出自莫斯科电子技术学院 1964—1966 年的入学考试。对考生错误的分析总是有教益的：它可以引出问题的不同侧面，点出某些细微之处，从而达到对基础知识更透彻的理解。

\pnote{原书成书于 1968/1973 年，"入学考试"指苏联工科高校的招生考试（口试占重要地位）。对今天的中国读者，本书的合适位置是：学完高中物理后的复习查漏与竞赛入门思维训练；它不是入门教材，作者在下文也是这样定位的。}

不过，这种分析做起来非常困难。正确的答案只有一个，错误的答案却可以有很多。任何问题都不可能预见全部错误答案——许多错误永远藏在口试学生令人着急的沉默背后。尽管如此，某些问题特有的错误答案总是反复出现；有许多问题几乎不可避免地会被答错。本书主要就以这类问题和习题为基础。

我们要提醒读者：本书绝不是一本涵盖大纲全部条目的教材。读者在书中找不到物理课程所要求的系统叙述。本书更像一篇随意讲出的故事，或者不如说，一场随意展开的讨论。因此，对那些想从头学习物理或系统整理物理知识的人，本书用处不大。它是为那些在考试前夕希望扩充物理知识的人准备的。

在我们心目中，理想的读者已经学完中学物理课程，对物理内容有清晰的整体图景，记得主要关系式，能引用各种定律，对所用的单位也相当熟悉。他正处于那种"悬置"状态：不再是中学生，也还不是正式的大学生——但他渴望成为大学生。如果为此需要扩展物理知识，本书可以帮他一把。

首先，我们希望本书能证明：死记硬背教材（哪怕是很好的教材）不仅令人厌倦，而且毫无成效。学生必须学会思考，学会琢磨材料本身，而不是单纯背诵。如果读者能在某种程度上获得这种认识，我们就认为自己的努力是值得的。

最后，我们感谢 G. 叶皮法诺夫教授，没有他的鼓励和宝贵帮助，本书不可能写成并出版。我们还衷心感谢莫斯科动力学院（Moscow Power Engineering Institute）物理系高级教师 E. N. 弗托罗夫（E.~N.~Vtorov）、副教授 A. G. 切尔托夫（A.~G.~Chertov）和 V. A. 法布里坎特教授（V.~A.~Fabrikant）对手稿提出的许多有益建议和建设性批评。

\begin{flushright}
\textit{L. 塔拉索夫}\\
\textit{A. 塔拉索娃}
\end{flushright}

\cleardoublepage
